\scanpage{615}\mspage{3}
% archive sheet number: 614

contient les iso, stable par chgt de base et par compositions
\unsure{transfinies|transfinie}. Si $\Phi$ est Q-saturé à gauche, il est stable par
\inserted{(petites)} sommes directes i.e. si \inserted{les}
$f_i:A_i\unsure{-|\to}B_i$ ($i\in I$, $I$ petit) sont dans $\Phi$,
$f=\coprod f_i:\coprod A_i\longrightarrow\coprod B_i$ aussi ;
et de $=$ pour \inserted{(petites)} sommes amalgamées finies ou
infinies (quand on se place dans un $A\backslash M$, au lieu de $M$).
On a l'énoncé dual pour $\Psi$, quand $\Psi$ est Q-saturé à droite.
\deleted{\unsure{En fait}}

En fait, il y a des notions plus faibles de Q-saturation que celle
que je viens de dire : Q-\uline{saturation simple} (à gauche :
$\Phi$ contient les iso, stable par composition binaire, par
cochgt de base), Q-\uline{saturation transfinie} (on exige de plus
stabilité par compositions transfinies), enfin Q-\uline{saturation
totale}, ou Q-saturation tout court\unsure{\S|)}, (quand on exige de
plus la stabilité par facteurs directs). On a les implications
\diagram{
\[
\left\{\vcenter{\hbox{
\begin{tikzcd}[ampersand replacement=\&,column sep=0pt,row sep=1em]
\text{Q-}\arrow[d,Rightarrow] \& |[anchor=west]|\text{\deleted{\unsure{sat}}clos à g}\\
\text{Q-}\arrow[d,Rightarrow] \& |[anchor=west]|\text{saturé (total)}\\
\text{Q-}\arrow[d,Rightarrow] \& |[anchor=west]|\text{saturé transfini}\\
\text{Q-} \& |[anchor=west]|\text{saturé simple}\, .
\end{tikzcd}
}}\right.\tag*{(1.6)}
\]
}

\deleted{\unsure{Un}} Pour $\Phi\subset\mathrm{Fl}\,M$, on a sa
Q-\uline{clôture} \inserted{(\uline{à gauche})} $\Phi^{\sim}$
(qui est la plus petite partie Q-close à gauche de
$\mathrm{Fl}\,M$ contenant $\Phi$), et de $=$ son Q-saturé
\inserted{(total)} à gauche, son Q-\uline{saturé transfini}, son
Q-\uline{saturé simple},
